% «Компьютер на 20 ваттах» — популярная версия («в дорогу»).
% Каждая глава paperback/NN.tex соответствует полной главе ../chapters/NN_*.tex с тем же номером;
% на сайте читатель переключает главу между версиями «Коротко» и «Полностью».
\documentclass[11pt,openany]{book}
\usepackage[paperwidth=13cm,paperheight=20cm,top=1.8cm,bottom=1.8cm,inner=1.6cm,outer=1.4cm]{geometry}
\usepackage[T2A]{fontenc}
\usepackage[utf8]{inputenc}
\usepackage[english,russian]{babel}
\usepackage{amsmath,amssymb}
\usepackage{xcolor}
\usepackage[unicode,colorlinks=true,linkcolor=blue!60!black]{hyperref}
\usepackage{microtype}
\input{paperback/pencil} % minimal colour-pencil illustrations, one per chapter
\setlength{\parskip}{0.3em}
\setlength{\emergencystretch}{2em} % narrow paperback page: allow a little extra stretch
\renewcommand{\chaptermark}[1]{\markboth{\small #1}{}}
% neuron type code: first four letters of the primary mediator (as in the full edition)
\newcommand{\nt}[1]{\textsf{\textbf{#1}}}
% pointer to the full edition at the end of each chapter
\newcommand{\fullversion}[1]{\par\bigskip\noindent\small\emph{Формулы, модели, задачи и источники --- в полной версии, глава~#1.}\normalsize}

\title{\Huge Компьютер на 20 ваттах\\[14pt] \Large Как вычисляет мозг --- коротко и просто}
\author{\Large К.~Кладько}
\date{}

\begin{document}
\frontmatter
\maketitle
\tableofcontents
\input{paperback/00_preface}
\mainmatter
\input{paperback/01}
\input{paperback/02}
\input{paperback/03}
\input{paperback/04}
\input{paperback/05}
\input{paperback/06}
\input{paperback/07}
\input{paperback/08}
\input{paperback/09}
\input{paperback/10}
\input{paperback/11}
\input{paperback/12}
\input{paperback/13}
\input{paperback/14}
\input{paperback/15}
\input{paperback/16}
\input{paperback/17}
\input{paperback/18}
\input{paperback/19}
\input{paperback/20}
\input{paperback/21}
\input{paperback/22}
\end{document}
